% Options for packages loaded elsewhere
% Options for packages loaded elsewhere
\PassOptionsToPackage{unicode}{hyperref}
\PassOptionsToPackage{hyphens}{url}
\PassOptionsToPackage{dvipsnames,svgnames,x11names}{xcolor}
%
\documentclass[
  11pt,
  letterpaper,
  DIV=11,
  numbers=noendperiod]{scrartcl}
\usepackage{xcolor}
\usepackage[margin=2cm]{geometry}
\usepackage{amsmath,amssymb}
\setcounter{secnumdepth}{5}
\usepackage{iftex}
\ifPDFTeX
  \usepackage[T1]{fontenc}
  \usepackage[utf8]{inputenc}
  \usepackage{textcomp} % provide euro and other symbols
\else % if luatex or xetex
  \usepackage{unicode-math} % this also loads fontspec
  \defaultfontfeatures{Scale=MatchLowercase}
  \defaultfontfeatures[\rmfamily]{Ligatures=TeX,Scale=1}
\fi
\usepackage{lmodern}
\ifPDFTeX\else
  % xetex/luatex font selection
\fi
% Use upquote if available, for straight quotes in verbatim environments
\IfFileExists{upquote.sty}{\usepackage{upquote}}{}
\IfFileExists{microtype.sty}{% use microtype if available
  \usepackage[]{microtype}
  \UseMicrotypeSet[protrusion]{basicmath} % disable protrusion for tt fonts
}{}
\makeatletter
\@ifundefined{KOMAClassName}{% if non-KOMA class
  \IfFileExists{parskip.sty}{%
    \usepackage{parskip}
  }{% else
    \setlength{\parindent}{0pt}
    \setlength{\parskip}{6pt plus 2pt minus 1pt}}
}{% if KOMA class
  \KOMAoptions{parskip=half}}
\makeatother
% Make \paragraph and \subparagraph free-standing
\makeatletter
\ifx\paragraph\undefined\else
  \let\oldparagraph\paragraph
  \renewcommand{\paragraph}{
    \@ifstar
      \xxxParagraphStar
      \xxxParagraphNoStar
  }
  \newcommand{\xxxParagraphStar}[1]{\oldparagraph*{#1}\mbox{}}
  \newcommand{\xxxParagraphNoStar}[1]{\oldparagraph{#1}\mbox{}}
\fi
\ifx\subparagraph\undefined\else
  \let\oldsubparagraph\subparagraph
  \renewcommand{\subparagraph}{
    \@ifstar
      \xxxSubParagraphStar
      \xxxSubParagraphNoStar
  }
  \newcommand{\xxxSubParagraphStar}[1]{\oldsubparagraph*{#1}\mbox{}}
  \newcommand{\xxxSubParagraphNoStar}[1]{\oldsubparagraph{#1}\mbox{}}
\fi
\makeatother

\usepackage{color}
\usepackage{fancyvrb}
\newcommand{\VerbBar}{|}
\newcommand{\VERB}{\Verb[commandchars=\\\{\}]}
\DefineVerbatimEnvironment{Highlighting}{Verbatim}{commandchars=\\\{\}}
% Add ',fontsize=\small' for more characters per line
\usepackage{framed}
\definecolor{shadecolor}{RGB}{241,243,245}
\newenvironment{Shaded}{\begin{snugshade}}{\end{snugshade}}
\newcommand{\AlertTok}[1]{\textcolor[rgb]{0.68,0.00,0.00}{#1}}
\newcommand{\AnnotationTok}[1]{\textcolor[rgb]{0.37,0.37,0.37}{#1}}
\newcommand{\AttributeTok}[1]{\textcolor[rgb]{0.40,0.45,0.13}{#1}}
\newcommand{\BaseNTok}[1]{\textcolor[rgb]{0.68,0.00,0.00}{#1}}
\newcommand{\BuiltInTok}[1]{\textcolor[rgb]{0.00,0.23,0.31}{#1}}
\newcommand{\CharTok}[1]{\textcolor[rgb]{0.13,0.47,0.30}{#1}}
\newcommand{\CommentTok}[1]{\textcolor[rgb]{0.37,0.37,0.37}{#1}}
\newcommand{\CommentVarTok}[1]{\textcolor[rgb]{0.37,0.37,0.37}{\textit{#1}}}
\newcommand{\ConstantTok}[1]{\textcolor[rgb]{0.56,0.35,0.01}{#1}}
\newcommand{\ControlFlowTok}[1]{\textcolor[rgb]{0.00,0.23,0.31}{\textbf{#1}}}
\newcommand{\DataTypeTok}[1]{\textcolor[rgb]{0.68,0.00,0.00}{#1}}
\newcommand{\DecValTok}[1]{\textcolor[rgb]{0.68,0.00,0.00}{#1}}
\newcommand{\DocumentationTok}[1]{\textcolor[rgb]{0.37,0.37,0.37}{\textit{#1}}}
\newcommand{\ErrorTok}[1]{\textcolor[rgb]{0.68,0.00,0.00}{#1}}
\newcommand{\ExtensionTok}[1]{\textcolor[rgb]{0.00,0.23,0.31}{#1}}
\newcommand{\FloatTok}[1]{\textcolor[rgb]{0.68,0.00,0.00}{#1}}
\newcommand{\FunctionTok}[1]{\textcolor[rgb]{0.28,0.35,0.67}{#1}}
\newcommand{\ImportTok}[1]{\textcolor[rgb]{0.00,0.46,0.62}{#1}}
\newcommand{\InformationTok}[1]{\textcolor[rgb]{0.37,0.37,0.37}{#1}}
\newcommand{\KeywordTok}[1]{\textcolor[rgb]{0.00,0.23,0.31}{\textbf{#1}}}
\newcommand{\NormalTok}[1]{\textcolor[rgb]{0.00,0.23,0.31}{#1}}
\newcommand{\OperatorTok}[1]{\textcolor[rgb]{0.37,0.37,0.37}{#1}}
\newcommand{\OtherTok}[1]{\textcolor[rgb]{0.00,0.23,0.31}{#1}}
\newcommand{\PreprocessorTok}[1]{\textcolor[rgb]{0.68,0.00,0.00}{#1}}
\newcommand{\RegionMarkerTok}[1]{\textcolor[rgb]{0.00,0.23,0.31}{#1}}
\newcommand{\SpecialCharTok}[1]{\textcolor[rgb]{0.37,0.37,0.37}{#1}}
\newcommand{\SpecialStringTok}[1]{\textcolor[rgb]{0.13,0.47,0.30}{#1}}
\newcommand{\StringTok}[1]{\textcolor[rgb]{0.13,0.47,0.30}{#1}}
\newcommand{\VariableTok}[1]{\textcolor[rgb]{0.07,0.07,0.07}{#1}}
\newcommand{\VerbatimStringTok}[1]{\textcolor[rgb]{0.13,0.47,0.30}{#1}}
\newcommand{\WarningTok}[1]{\textcolor[rgb]{0.37,0.37,0.37}{\textit{#1}}}

\usepackage{longtable,booktabs,array}
\newcounter{none} % for unnumbered tables
\usepackage{calc} % for calculating minipage widths
% Correct order of tables after \paragraph or \subparagraph
\usepackage{etoolbox}
\makeatletter
\patchcmd\longtable{\par}{\if@noskipsec\mbox{}\fi\par}{}{}
\makeatother
% Allow footnotes in longtable head/foot
\IfFileExists{footnotehyper.sty}{\usepackage{footnotehyper}}{\usepackage{footnote}}
\makesavenoteenv{longtable}
\usepackage{graphicx}
\makeatletter
\newsavebox\pandoc@box
\newcommand*\pandocbounded[1]{% scales image to fit in text height/width
  \sbox\pandoc@box{#1}%
  \Gscale@div\@tempa{\textheight}{\dimexpr\ht\pandoc@box+\dp\pandoc@box\relax}%
  \Gscale@div\@tempb{\linewidth}{\wd\pandoc@box}%
  \ifdim\@tempb\p@<\@tempa\p@\let\@tempa\@tempb\fi% select the smaller of both
  \ifdim\@tempa\p@<\p@\scalebox{\@tempa}{\usebox\pandoc@box}%
  \else\usebox{\pandoc@box}%
  \fi%
}
% Set default figure placement to htbp
\def\fps@figure{htbp}
\makeatother





\setlength{\emergencystretch}{3em} % prevent overfull lines

\providecommand{\tightlist}{%
  \setlength{\itemsep}{0pt}\setlength{\parskip}{0pt}}



 


\KOMAoption{captions}{tableheading}
\makeatletter
\@ifpackageloaded{tcolorbox}{}{\usepackage[skins,breakable]{tcolorbox}}
\@ifpackageloaded{fontawesome5}{}{\usepackage{fontawesome5}}
\definecolor{quarto-callout-color}{HTML}{909090}
\definecolor{quarto-callout-note-color}{HTML}{0758E5}
\definecolor{quarto-callout-important-color}{HTML}{CC1914}
\definecolor{quarto-callout-warning-color}{HTML}{EB9113}
\definecolor{quarto-callout-tip-color}{HTML}{00A047}
\definecolor{quarto-callout-caution-color}{HTML}{FC5300}
\definecolor{quarto-callout-color-frame}{HTML}{acacac}
\definecolor{quarto-callout-note-color-frame}{HTML}{4582ec}
\definecolor{quarto-callout-important-color-frame}{HTML}{d9534f}
\definecolor{quarto-callout-warning-color-frame}{HTML}{f0ad4e}
\definecolor{quarto-callout-tip-color-frame}{HTML}{02b875}
\definecolor{quarto-callout-caution-color-frame}{HTML}{fd7e14}
\makeatother
\makeatletter
\@ifpackageloaded{caption}{}{\usepackage{caption}}
\AtBeginDocument{%
\ifdefined\contentsname
  \renewcommand*\contentsname{Table of contents}
\else
  \newcommand\contentsname{Table of contents}
\fi
\ifdefined\listfigurename
  \renewcommand*\listfigurename{List of Figures}
\else
  \newcommand\listfigurename{List of Figures}
\fi
\ifdefined\listtablename
  \renewcommand*\listtablename{List of Tables}
\else
  \newcommand\listtablename{List of Tables}
\fi
\ifdefined\figurename
  \renewcommand*\figurename{Figure}
\else
  \newcommand\figurename{Figure}
\fi
\ifdefined\tablename
  \renewcommand*\tablename{Table}
\else
  \newcommand\tablename{Table}
\fi
}
\@ifpackageloaded{float}{}{\usepackage{float}}
\floatstyle{ruled}
\@ifundefined{c@chapter}{\newfloat{codelisting}{h}{lop}}{\newfloat{codelisting}{h}{lop}[chapter]}
\floatname{codelisting}{Listing}
\newcommand*\listoflistings{\listof{codelisting}{List of Listings}}
\makeatother
\makeatletter
\makeatother
\makeatletter
\@ifpackageloaded{caption}{}{\usepackage{caption}}
\@ifpackageloaded{subcaption}{}{\usepackage{subcaption}}
\makeatother
\usepackage{bookmark}
\IfFileExists{xurl.sty}{\usepackage{xurl}}{} % add URL line breaks if available
\urlstyle{same}
\hypersetup{
  pdftitle={Representações Numéricas no Computador},
  pdfauthor={Prof.~Dr.~Raphael Teixeira},
  colorlinks=true,
  linkcolor={blue},
  filecolor={Maroon},
  citecolor={Blue},
  urlcolor={Blue},
  pdfcreator={LaTeX via pandoc}}


\title{Representações Numéricas no Computador}
\usepackage{etoolbox}
\makeatletter
\providecommand{\subtitle}[1]{% add subtitle to \maketitle
  \apptocmd{\@title}{\par {\large #1 \par}}{}{}
}
\makeatother
\subtitle{Fundamentos Computacionais para Cálculo Numérico}
\author{Prof.~Dr.~Raphael Teixeira}
\date{2026-08-15}
\begin{document}
\maketitle

\renewcommand*\contentsname{Table of contents}
{
\setcounter{tocdepth}{3}
\tableofcontents
}

\section{Introdução: O abismo entre matemática e
computação}\label{introduuxe7uxe3o-o-abismo-entre-matemuxe1tica-e-computauxe7uxe3o}

Na matemática, operamos com números reais em \textbf{precisão infinita}.
O conjunto \(\mathbb{R}\) é denso: entre dois reais quaisquer existe
sempre outro real. Podemos trabalhar com \(\pi\), \(e\), ou \(\sqrt{2}\)
sem restrições.

\textbf{No computador, tudo muda.}

Um computador tem memória \textbf{finita}. Se cada número real
precisasse de infinitos dígitos para representação exata, seria
impossível armazená-lo. Logo, \textbf{todo número é uma aproximação}.

\subsection{Por que isso importa para Cálculo
Numérico}\label{por-que-isso-importa-para-cuxe1lculo-numuxe9rico}

\begin{enumerate}
\def\labelenumi{\arabic{enumi}.}
\tightlist
\item
  \textbf{Erros de arredondamento} --- cada operação introduz pequenos
  erros
\item
  \textbf{Propagação de erros} --- erros crescem ao longo de
  algoritmos\\
\item
  \textbf{Instabilidade numérica} --- alguns algoritmos amplificam
  erros, outros os controlam
\item
  \textbf{Escolha de método} --- um algoritmo ``bom'' matematicamente
  pode ser ruim em máquina
\end{enumerate}

\begin{tcolorbox}[enhanced jigsaw, arc=.35mm, bottomrule=.15mm, bottomtitle=1mm, breakable, colback=white, colbacktitle=quarto-callout-warning-color!10!white, colframe=quarto-callout-warning-color-frame, coltitle=black, left=2mm, leftrule=.75mm, opacityback=0, opacitybacktitle=0.6, rightrule=.15mm, title=\textcolor{quarto-callout-warning-color}{\faExclamationTriangle}\hspace{0.5em}{Exemplo clássico}, titlerule=0mm, toprule=.15mm, toptitle=1mm]

Resolver \(x^2 - 1000x + 1 = 0\) com a fórmula de Bhaskara pode resultar
em \textbf{catastrófica perda de dígitos significativos} devido a
cancelamento numérico.

\end{tcolorbox}

\begin{center}\rule{0.5\linewidth}{0.5pt}\end{center}

\section{Sistemas de Numeração
Posicional}\label{sistemas-de-numerauxe7uxe3o-posicional}

\subsection{Base decimal (base 10)}\label{base-decimal-base-10}

Um número em base 10 é representado como:
\[N = d_n \cdot 10^n + d_{n-1} \cdot 10^{n-1} + \cdots + d_1 \cdot 10 + d_0 + d_{-1} \cdot 10^{-1} + \cdots\]

onde cada \(d_i \in \{0, 1, \ldots, 9\}\).

\textbf{Exemplo:}
\(1234.56 = 1 \cdot 10^3 + 2 \cdot 10^2 + 3 \cdot 10 + 4 + 5 \cdot 10^{-1} + 6 \cdot 10^{-2}\)

\subsection{Base binária (base 2)}\label{base-binuxe1ria-base-2}

Um número em base 2 usa apenas dígitos \(\{0, 1\}\) (bits):
\[N = b_n \cdot 2^n + b_{n-1} \cdot 2^{n-1} + \cdots + b_0 + b_{-1} \cdot 2^{-1} + b_{-2} \cdot 2^{-2} + \cdots\]

onde cada \(b_i \in \{0, 1\}\).

\begin{tcolorbox}[enhanced jigsaw, arc=.35mm, bottomrule=.15mm, bottomtitle=1mm, breakable, colback=white, colbacktitle=quarto-callout-tip-color!10!white, colframe=quarto-callout-tip-color-frame, coltitle=black, left=2mm, leftrule=.75mm, opacityback=0, opacitybacktitle=0.6, rightrule=.15mm, title=\textcolor{quarto-callout-tip-color}{\faLightbulb}\hspace{0.5em}{Exemplo}, titlerule=0mm, toprule=.15mm, toptitle=1mm]

\((1101.101)_2 = 1 \cdot 2^3 + 1 \cdot 2^2 + 0 \cdot 2^1 + 1 \cdot 2^0 + 1 \cdot 2^{-1} + 0 \cdot 2^{-2} + 1 \cdot 2^{-3}\)

\(= 8 + 4 + 0 + 1 + 0.5 + 0 + 0.125 = 13.625\) em decimal

\end{tcolorbox}

\subsection{Conversão decimal para
binário}\label{conversuxe3o-decimal-para-binuxe1rio}

\textbf{Parte inteira:} divisões sucessivas por 2, tomando os restos.

\textbf{Parte fracionária:} multiplicações sucessivas por 2, tomando as
partes inteiras.

\begin{tcolorbox}[enhanced jigsaw, arc=.35mm, bottomrule=.15mm, bottomtitle=1mm, breakable, colback=white, colbacktitle=quarto-callout-tip-color!10!white, colframe=quarto-callout-tip-color-frame, coltitle=black, left=2mm, leftrule=.75mm, opacityback=0, opacitybacktitle=0.6, rightrule=.15mm, title=\textcolor{quarto-callout-tip-color}{\faLightbulb}\hspace{0.5em}{Exemplo: Converter 0.375}, titlerule=0mm, toprule=.15mm, toptitle=1mm]

\begin{verbatim}
0.375 × 2 = 0.75    (dígito: 0)
0.75 × 2 = 1.5      (dígito: 1)
0.5 × 2 = 1.0       (dígito: 1)
\end{verbatim}

Logo, \((0.375)_{10} = (0.011)_2\)

\end{tcolorbox}

\begin{tcolorbox}[enhanced jigsaw, arc=.35mm, bottomrule=.15mm, bottomtitle=1mm, breakable, colback=white, colbacktitle=quarto-callout-warning-color!10!white, colframe=quarto-callout-warning-color-frame, coltitle=black, left=2mm, leftrule=.75mm, opacityback=0, opacitybacktitle=0.6, rightrule=.15mm, title=\textcolor{quarto-callout-warning-color}{\faExclamationTriangle}\hspace{0.5em}{Atenção}, titlerule=0mm, toprule=.15mm, toptitle=1mm]

Nem todo número decimal tem representação binária \textbf{finita}!

Por exemplo, \((0.1)_{10} = (0.0001100110011\ldots)_2\) é dízima
periódica infinita.

\textbf{Implicação:} 0.1 não pode ser representado exatamente em
computador!

\end{tcolorbox}

\begin{Shaded}
\begin{Highlighting}[]
\CommentTok{\# Demonstração: 0.1 não é exato em binário}
\NormalTok{x }\OperatorTok{=} \FloatTok{0.1}
\BuiltInTok{print}\NormalTok{(}\SpecialStringTok{f"0.1 em Python (16 casas): }\SpecialCharTok{\{}\NormalTok{x}\SpecialCharTok{:.16f\}}\SpecialStringTok{"}\NormalTok{)}
\BuiltInTok{print}\NormalTok{(}\SpecialStringTok{f"0.1 em Python (20 casas): }\SpecialCharTok{\{}\NormalTok{x}\SpecialCharTok{:.20f\}}\SpecialStringTok{"}\NormalTok{)}

\CommentTok{\# Isso causa problemas}
\BuiltInTok{print}\NormalTok{(}\SpecialStringTok{f"}\CharTok{\textbackslash{}n}\SpecialStringTok{0.1 + 0.2 = }\SpecialCharTok{\{}\FloatTok{0.1} \OperatorTok{+} \FloatTok{0.2}\SpecialCharTok{:.20f\}}\SpecialStringTok{"}\NormalTok{)}
\BuiltInTok{print}\NormalTok{(}\SpecialStringTok{f"Esperado:  0.3"}\NormalTok{)}
\end{Highlighting}
\end{Shaded}

\begin{verbatim}
0.1 em Python (16 casas): 0.1000000000000000
0.1 em Python (20 casas): 0.10000000000000000555

0.1 + 0.2 = 0.30000000000000004441
Esperado:  0.3
\end{verbatim}

\begin{center}\rule{0.5\linewidth}{0.5pt}\end{center}

\section{Inteiros em Máquina: Aritmética de Ponto
Fixo}\label{inteiros-em-muxe1quina-aritmuxe9tica-de-ponto-fixo}

\subsection{Inteiros sem sinal (unsigned
integers)}\label{inteiros-sem-sinal-unsigned-integers}

Um inteiro sem sinal com \(n\) bits pode representar inteiros de \(0\) a
\(2^n - 1\).

{\def\LTcaptype{none} % do not increment counter
\begin{longtable}[]{@{}lll@{}}
\toprule\noalign{}
Bits & Intervalo & Uso \\
\midrule\noalign{}
\endhead
\bottomrule\noalign{}
\endlastfoot
8 & {[}0, 255{]} & 1 byte \\
16 & {[}0, 65535{]} & 2 bytes \\
32 & {[}0, 4.3 × 10⁹{]} & 4 bytes \\
64 & {[}0, 1.8 × 10¹⁹{]} & 8 bytes \\
\end{longtable}
}

\subsection{Inteiros com sinal: Complemento de
dois}\label{inteiros-com-sinal-complemento-de-dois}

\textbf{Padrão moderno} para representar negativos (sem ambiguidade).

Um inteiro com sinal em \(n\) bits é representado por:
\[N = -b_{n-1} \cdot 2^{n-1} + \sum_{i=0}^{n-2} b_i \cdot 2^i\]

onde o bit mais significativo (MSB) é o \textbf{bit de sinal}.

\textbf{Intervalo:} \(-2^{n-1}\) a \(2^{n-1} - 1\)

\begin{tcolorbox}[enhanced jigsaw, arc=.35mm, bottomrule=.15mm, bottomtitle=1mm, breakable, colback=white, colbacktitle=quarto-callout-tip-color!10!white, colframe=quarto-callout-tip-color-frame, coltitle=black, left=2mm, leftrule=.75mm, opacityback=0, opacitybacktitle=0.6, rightrule=.15mm, title=\textcolor{quarto-callout-tip-color}{\faLightbulb}\hspace{0.5em}{Exemplo com 8 bits}, titlerule=0mm, toprule=.15mm, toptitle=1mm]

\begin{itemize}
\tightlist
\item
  \((00000101)_2 = 5\)
\item
  \((11111011)_2 = -5\)
\item
  \((10000000)_2 = -128\) (menor valor)
\item
  \((01111111)_2 = 127\) (maior valor)
\end{itemize}

\textbf{Intervalo:} \([-128, 127]\)

\end{tcolorbox}

\subsection{Overflow inteiro}\label{overflow-inteiro}

Se o resultado de uma operação inteira exceder o intervalo
representável, ocorre \textbf{overflow}.

\begin{tcolorbox}[enhanced jigsaw, arc=.35mm, bottomrule=.15mm, bottomtitle=1mm, breakable, colback=white, colbacktitle=quarto-callout-warning-color!10!white, colframe=quarto-callout-warning-color-frame, coltitle=black, left=2mm, leftrule=.75mm, opacityback=0, opacitybacktitle=0.6, rightrule=.15mm, title=\textcolor{quarto-callout-warning-color}{\faExclamationTriangle}\hspace{0.5em}{Comportamentos diferentes por linguagem}, titlerule=0mm, toprule=.15mm, toptitle=1mm]

\begin{itemize}
\tightlist
\item
  \textbf{C/Fortran:} comportamento indefinido ou truncamento silencioso
\item
  \textbf{Python:} inteiros crescem dinamicamente (sem overflow)
\item
  \textbf{Java:} wraparound (volta ao intervalo válido)
\end{itemize}

\textbf{Cuidado:} Em linguagens compiladas, overflow silencioso é
perigoso!

\end{tcolorbox}

\begin{center}\rule{0.5\linewidth}{0.5pt}\end{center}

\section{Números Reais em Máquina: Ponto
Flutuante}\label{nuxfameros-reais-em-muxe1quina-ponto-flutuante}

\subsection{Notação científica}\label{notauxe7uxe3o-cientuxedfica}

Em matemática, qualquer número real pode ser escrito em notação
científica: \[x = \pm d_0.d_1 d_2 d_3 \ldots \times 10^e\]

onde \(d_0 \in \{1, 2, \ldots, 9\}\) (um dígito antes do ponto) e \(e\)
é o expoente.

\textbf{Exemplo:} \(1234.56 = 1.23456 \times 10^3\)

\subsection{Representação de ponto flutuante
normalizada}\label{representauxe7uxe3o-de-ponto-flutuante-normalizada}

Em base 2, um número de ponto flutuante tem a forma:
\[x = \pm (1.d_1 d_2 d_3 \ldots) \times 2^e\]

onde: - O bit implícito na posição \(2^0\) é sempre 1 (em representação
normalizada) - \(d_i \in \{0, 1\}\) são os dígitos da mantissa (ou
fração) - \(e\) é o expoente

\textbf{Componentes:} 1. \textbf{Sinal:} 1 bit (0 = positivo, 1 =
negativo) 2. \textbf{Expoente:} \(e\) bits (codificado com bias) 3.
\textbf{Mantissa/Fração:} \(p\) bits (precisão)

\textbf{Precisão total:} \(p + 1\) dígitos binários (o ``+1'' é o bit
implícito)

\begin{center}\rule{0.5\linewidth}{0.5pt}\end{center}

\section{Padrão IEEE 754}\label{padruxe3o-ieee-754}

IEEE 754 (IEEE Standard for Floating-Point Arithmetic) é o padrão
internacional para aritmética de ponto flutuante. Garante consistência
entre plataformas.

\subsection{Formatos principais}\label{formatos-principais}

{\def\LTcaptype{none} % do not increment counter
\begin{longtable}[]{@{}lllllll@{}}
\toprule\noalign{}
Tipo & Bits & Sinal & Exp & Mant & Bias & Uso \\
\midrule\noalign{}
\endhead
\bottomrule\noalign{}
\endlastfoot
\textbf{float32} & 32 & 1 & 8 & 23 & 127 & Gráficos, GPUs \\
\textbf{float64} & 64 & 1 & 11 & 52 & 1023 & Ciência, padrão Python \\
\textbf{float80} & 80 & 1 & 15 & 64 & 16383 & Cálculos intermediários \\
\end{longtable}
}

\begin{tcolorbox}[enhanced jigsaw, arc=.35mm, bottomrule=.15mm, bottomtitle=1mm, breakable, colback=white, colbacktitle=quarto-callout-note-color!10!white, colframe=quarto-callout-note-color-frame, coltitle=black, left=2mm, leftrule=.75mm, opacityback=0, opacitybacktitle=0.6, rightrule=.15mm, title=\textcolor{quarto-callout-note-color}{\faInfo}\hspace{0.5em}{Foco neste curso}, titlerule=0mm, toprule=.15mm, toptitle=1mm]

Trabalharemos principalmente com \textbf{float64} (precisão dupla, 64
bits), que é o padrão em Python, MATLAB e linguagens científicas.

\end{tcolorbox}

\subsection{Estrutura de um número double (64
bits)}\label{estrutura-de-um-nuxfamero-double-64-bits}

\begin{verbatim}
┌─────────┬───────────────────┬────────────────────────────────────────────────────┐
│  Sinal  │   Expoente (11)   │              Mantissa (52 bits)                    │
│ (1 bit) │                   │                                                    │
└─────────┴───────────────────┴────────────────────────────────────────────────────┘
\end{verbatim}

\textbf{Decodificação:}
\[\text{valor} = (-1)^S \times 2^{E - 1023} \times (1.M)_2\]

onde \((1.M)_2 = 1 + \sum_{i=1}^{52} M_i \cdot 2^{-i}\)

\begin{tcolorbox}[enhanced jigsaw, arc=.35mm, bottomrule=.15mm, bottomtitle=1mm, breakable, colback=white, colbacktitle=quarto-callout-tip-color!10!white, colframe=quarto-callout-tip-color-frame, coltitle=black, left=2mm, leftrule=.75mm, opacityback=0, opacitybacktitle=0.6, rightrule=.15mm, title=\textcolor{quarto-callout-tip-color}{\faLightbulb}\hspace{0.5em}{Exemplo de decodificação}, titlerule=0mm, toprule=.15mm, toptitle=1mm]

Considere o padrão IEEE 754 double:

\begin{verbatim}
0 01111111100 0000000000000000000000000000000000000000000000000000
\end{verbatim}

\textbf{Decodificação:} - \(S = 0\) (positivo) -
\(E = (01111111100)_2 = 1020\) (expoente com bias) - \(M = 0\) (mantissa
nula) - Expoente real: \(1020 - 1023 = -3\) - \textbf{Valor:}
\(1 \times 2^{-3} \times 1.0 = 0.125\)

\end{tcolorbox}

\subsection{Valores especiais em IEEE
754}\label{valores-especiais-em-ieee-754}

{\def\LTcaptype{none} % do not increment counter
\begin{longtable}[]{@{}lllll@{}}
\toprule\noalign{}
Padrão & Sinal & Expoente & Mantissa & Significado \\
\midrule\noalign{}
\endhead
\bottomrule\noalign{}
\endlastfoot
Zeros & ±0 & 0 & 0 & +0, -0 \\
Desnormalizados & ± & 0 & ≠ 0 & Muito pequenos (precisão reduzida) \\
Normalizados & ± & 1 a 2046 & qualquer & Números regulares \\
Infinitos & ± & máximo & 0 & \(\pm \infty\) \\
NaN & qualquer & máximo & ≠ 0 & Not a Number \\
\end{longtable}
}

\begin{tcolorbox}[enhanced jigsaw, arc=.35mm, bottomrule=.15mm, bottomtitle=1mm, breakable, colback=white, colbacktitle=quarto-callout-tip-color!10!white, colframe=quarto-callout-tip-color-frame, coltitle=black, left=2mm, leftrule=.75mm, opacityback=0, opacitybacktitle=0.6, rightrule=.15mm, title=\textcolor{quarto-callout-tip-color}{\faLightbulb}\hspace{0.5em}{Exemplos de operações especiais}, titlerule=0mm, toprule=.15mm, toptitle=1mm]

\begin{Shaded}
\begin{Highlighting}[]
\BuiltInTok{print}\NormalTok{(}\DecValTok{1}\OperatorTok{/}\DecValTok{0}\NormalTok{)          }\CommentTok{\# inf}
\BuiltInTok{print}\NormalTok{(}\OperatorTok{{-}}\DecValTok{1}\OperatorTok{/}\DecValTok{0}\NormalTok{)         }\CommentTok{\# {-}inf}
\BuiltInTok{print}\NormalTok{(}\DecValTok{0}\OperatorTok{/}\DecValTok{0}\NormalTok{)          }\CommentTok{\# nan}
\BuiltInTok{print}\NormalTok{(}\BuiltInTok{float}\NormalTok{(}\StringTok{\textquotesingle{}inf\textquotesingle{}}\NormalTok{) }\OperatorTok{{-}} \BuiltInTok{float}\NormalTok{(}\StringTok{\textquotesingle{}inf\textquotesingle{}}\NormalTok{))  }\CommentTok{\# nan}
\end{Highlighting}
\end{Shaded}

\end{tcolorbox}

\subsection{Números desnormalizados
(subnormal)}\label{nuxfameros-desnormalizados-subnormal}

Quando o expoente atinge seu mínimo (\(E = 0\)), números menores que o
``mínimo normal'' podem ser representados com precisão reduzida:

\[x = \pm (0.d_1 d_2 \ldots) \times 2^{e_{\text{min}}}\]

Isto fornece \textbf{underflow gradual} ao invés de queda abrupta para
zero.

\begin{center}\rule{0.5\linewidth}{0.5pt}\end{center}

\section{Precisão e Erros de
Máquina}\label{precisuxe3o-e-erros-de-muxe1quina}

\subsection{\texorpdfstring{Machine epsilon
(\(\epsilon\))}{Machine epsilon (\textbackslash epsilon)}}\label{machine-epsilon-epsilon}

\textbf{Definição:} Menor número \(\epsilon\) tal que
\(1 + \epsilon \neq 1\) em aritmética de ponto flutuante.

\textbf{Fórmula:} \[\epsilon = 2^{-p}\]

onde \(p\) é o número de dígitos binários de precisão (após o bit
implícito).

{\def\LTcaptype{none} % do not increment counter
\begin{longtable}[]{@{}llll@{}}
\toprule\noalign{}
Formato & \(p\) & Machine epsilon & Valor aproximado \\
\midrule\noalign{}
\endhead
\bottomrule\noalign{}
\endlastfoot
float32 & 24 & \(2^{-23}\) & \(1.19 \times 10^{-7}\) \\
float64 & 53 & \(2^{-52}\) & \(2.22 \times 10^{-16}\) \\
\end{longtable}
}

\begin{tcolorbox}[enhanced jigsaw, arc=.35mm, bottomrule=.15mm, bottomtitle=1mm, breakable, colback=white, colbacktitle=quarto-callout-warning-color!10!white, colframe=quarto-callout-warning-color-frame, coltitle=black, left=2mm, leftrule=.75mm, opacityback=0, opacitybacktitle=0.6, rightrule=.15mm, title=\textcolor{quarto-callout-warning-color}{\faExclamationTriangle}\hspace{0.5em}{Interpretação}, titlerule=0mm, toprule=.15mm, toptitle=1mm]

Machine epsilon é uma medida da \textbf{precisão relativa}: em máquina,
qualquer número é representado com erro relativo no máximo da ordem de
\(\epsilon / 2\).

\end{tcolorbox}

\begin{Shaded}
\begin{Highlighting}[]
\CommentTok{\# Verificar machine epsilon em Python}
\ImportTok{import}\NormalTok{ numpy }\ImportTok{as}\NormalTok{ np}

\NormalTok{eps32 }\OperatorTok{=}\NormalTok{ np.finfo(np.float32).eps}
\NormalTok{eps64 }\OperatorTok{=}\NormalTok{ np.finfo(np.float64).eps}

\BuiltInTok{print}\NormalTok{(}\SpecialStringTok{f"Machine epsilon (float32): }\SpecialCharTok{\{}\NormalTok{eps32}\SpecialCharTok{\}}\SpecialStringTok{"}\NormalTok{)}
\BuiltInTok{print}\NormalTok{(}\SpecialStringTok{f"Machine epsilon (float64): }\SpecialCharTok{\{}\NormalTok{eps64}\SpecialCharTok{\}}\SpecialStringTok{"}\NormalTok{)}

\BuiltInTok{print}\NormalTok{(}\SpecialStringTok{f"}\CharTok{\textbackslash{}n}\SpecialStringTok{Teórico (2\^{}{-}23): }\SpecialCharTok{\{}\DecValTok{2}\OperatorTok{**}\NormalTok{(}\OperatorTok{{-}}\DecValTok{23}\NormalTok{)}\SpecialCharTok{\}}\SpecialStringTok{"}\NormalTok{)}
\BuiltInTok{print}\NormalTok{(}\SpecialStringTok{f"Teórico (2\^{}{-}52): }\SpecialCharTok{\{}\DecValTok{2}\OperatorTok{**}\NormalTok{(}\OperatorTok{{-}}\DecValTok{52}\NormalTok{)}\SpecialCharTok{\}}\SpecialStringTok{"}\NormalTok{)}

\CommentTok{\# Verificar experimentalmente}
\KeywordTok{def}\NormalTok{ find\_epsilon(dtype):}
\NormalTok{    eps }\OperatorTok{=}\NormalTok{ dtype(}\FloatTok{1.0}\NormalTok{)}
    \ControlFlowTok{while}\NormalTok{ dtype(}\FloatTok{1.0}\NormalTok{) }\OperatorTok{+}\NormalTok{ eps }\OperatorTok{!=}\NormalTok{ dtype(}\FloatTok{1.0}\NormalTok{):}
\NormalTok{        eps }\OperatorTok{=}\NormalTok{ eps }\OperatorTok{/}\NormalTok{ dtype(}\FloatTok{2.0}\NormalTok{)}
    \ControlFlowTok{return}\NormalTok{ eps }\OperatorTok{*} \DecValTok{2}  \CommentTok{\# Volta um passo}

\BuiltInTok{print}\NormalTok{(}\SpecialStringTok{f"}\CharTok{\textbackslash{}n}\SpecialStringTok{Experimental float64: }\SpecialCharTok{\{}\NormalTok{find\_epsilon(np.float64)}\SpecialCharTok{\}}\SpecialStringTok{"}\NormalTok{)}
\end{Highlighting}
\end{Shaded}

\begin{verbatim}
Machine epsilon (float32): 1.1920928955078125e-07
Machine epsilon (float64): 2.220446049250313e-16

Teórico (2^-23): 1.1920928955078125e-07
Teórico (2^-52): 2.220446049250313e-16

Experimental float64: 2.220446049250313e-16
\end{verbatim}

\subsection{Maior e menor número
representável}\label{maior-e-menor-nuxfamero-representuxe1vel}

{\def\LTcaptype{none} % do not increment counter
\begin{longtable}[]{@{}
  >{\raggedright\arraybackslash}p{(\linewidth - 4\tabcolsep) * \real{0.2368}}
  >{\raggedright\arraybackslash}p{(\linewidth - 4\tabcolsep) * \real{0.5789}}
  >{\raggedright\arraybackslash}p{(\linewidth - 4\tabcolsep) * \real{0.1842}}@{}}
\toprule\noalign{}
\begin{minipage}[b]{\linewidth}\raggedright
Formato
\end{minipage} & \begin{minipage}[b]{\linewidth}\raggedright
Menor positivo normal
\end{minipage} & \begin{minipage}[b]{\linewidth}\raggedright
Maior
\end{minipage} \\
\midrule\noalign{}
\endhead
\bottomrule\noalign{}
\endlastfoot
float32 & \(\approx 1.18 \times 10^{-38}\) &
\(\approx 3.40 \times 10^{38}\) \\
float64 & \(\approx 2.23 \times 10^{-308}\) &
\(\approx 1.80 \times 10^{308}\) \\
\end{longtable}
}

\begin{Shaded}
\begin{Highlighting}[]
\CommentTok{\# Verificar limites em Python}
\NormalTok{info32 }\OperatorTok{=}\NormalTok{ np.finfo(np.float32)}
\NormalTok{info64 }\OperatorTok{=}\NormalTok{ np.finfo(np.float64)}

\BuiltInTok{print}\NormalTok{(}\StringTok{"Float32:"}\NormalTok{)}
\BuiltInTok{print}\NormalTok{(}\SpecialStringTok{f"  Menor: }\SpecialCharTok{\{}\NormalTok{info32}\SpecialCharTok{.}\NormalTok{tiny}\SpecialCharTok{\}}\SpecialStringTok{"}\NormalTok{)}
\BuiltInTok{print}\NormalTok{(}\SpecialStringTok{f"  Maior: }\SpecialCharTok{\{}\NormalTok{info32}\SpecialCharTok{.}\BuiltInTok{max}\SpecialCharTok{\}}\SpecialStringTok{"}\NormalTok{)}

\BuiltInTok{print}\NormalTok{(}\StringTok{"}\CharTok{\textbackslash{}n}\StringTok{Float64:"}\NormalTok{)}
\BuiltInTok{print}\NormalTok{(}\SpecialStringTok{f"  Menor: }\SpecialCharTok{\{}\NormalTok{info64}\SpecialCharTok{.}\NormalTok{tiny}\SpecialCharTok{\}}\SpecialStringTok{"}\NormalTok{)}
\BuiltInTok{print}\NormalTok{(}\SpecialStringTok{f"  Maior: }\SpecialCharTok{\{}\NormalTok{info64}\SpecialCharTok{.}\BuiltInTok{max}\SpecialCharTok{\}}\SpecialStringTok{"}\NormalTok{)}
\end{Highlighting}
\end{Shaded}

\begin{verbatim}
Float32:
  Menor: 1.1754943508222875e-38
  Maior: 3.4028234663852886e+38

Float64:
  Menor: 2.2250738585072014e-308
  Maior: 1.7976931348623157e+308
\end{verbatim}

\subsection{Dígitos significativos
decimais}\label{duxedgitos-significativos-decimais}

Número de dígitos decimais significativos:
\[\text{# dígitos decimais} \approx p \times \log_{10}(2)\]

\begin{itemize}
\tightlist
\item
  \textbf{float32} (\(p = 24\)): \(\approx 24 \times 0.301 \approx 7\)
  dígitos
\item
  \textbf{float64} (\(p = 53\)): \(\approx 53 \times 0.301 \approx 16\)
  dígitos
\end{itemize}

\begin{Shaded}
\begin{Highlighting}[]
\CommentTok{\# Demonstração}
\NormalTok{x\_float32 }\OperatorTok{=}\NormalTok{ np.float32(}\FloatTok{1.23456789}\NormalTok{)}
\NormalTok{x\_float64 }\OperatorTok{=}\NormalTok{ np.float64(}\FloatTok{1.23456789}\NormalTok{)}

\BuiltInTok{print}\NormalTok{(}\SpecialStringTok{f"Entrada:  1.23456789"}\NormalTok{)}
\BuiltInTok{print}\NormalTok{(}\SpecialStringTok{f"float32: }\SpecialCharTok{\{}\NormalTok{x\_float32}\SpecialCharTok{:.8f\}}\SpecialStringTok{"}\NormalTok{)}
\BuiltInTok{print}\NormalTok{(}\SpecialStringTok{f"float64: }\SpecialCharTok{\{}\NormalTok{x\_float64}\SpecialCharTok{:.16f\}}\SpecialStringTok{"}\NormalTok{)}

\CommentTok{\# Valores maiores}
\NormalTok{x\_large32 }\OperatorTok{=}\NormalTok{ np.float32(}\FloatTok{123456789.0}\NormalTok{)}
\NormalTok{x\_large64 }\OperatorTok{=}\NormalTok{ np.float64(}\FloatTok{123456789.0}\NormalTok{)}

\BuiltInTok{print}\NormalTok{(}\SpecialStringTok{f"}\CharTok{\textbackslash{}n}\SpecialStringTok{Entrada:  123456789.0"}\NormalTok{)}
\BuiltInTok{print}\NormalTok{(}\SpecialStringTok{f"float32: }\SpecialCharTok{\{}\NormalTok{x\_large32}\SpecialCharTok{:.1f\}}\SpecialStringTok{"}\NormalTok{)}
\BuiltInTok{print}\NormalTok{(}\SpecialStringTok{f"float64: }\SpecialCharTok{\{}\NormalTok{x\_large64}\SpecialCharTok{:.10f\}}\SpecialStringTok{"}\NormalTok{)}
\end{Highlighting}
\end{Shaded}

\begin{verbatim}
Entrada:  1.23456789
float32: 1.23456788
float64: 1.2345678899999999

Entrada:  123456789.0
float32: 123456792.0
float64: 123456789.0000000000
\end{verbatim}

\begin{center}\rule{0.5\linewidth}{0.5pt}\end{center}

\section{Arredondamento}\label{arredondamento}

\subsection{Estratégias de
arredondamento}\label{estratuxe9gias-de-arredondamento}

Quando um número não cabe exatamente em \(p\) bits, deve-se arredondar.
IEEE 754 suporta 4 modos:

\begin{enumerate}
\def\labelenumi{\arabic{enumi}.}
\tightlist
\item
  \textbf{Round to nearest (padrão):} arredonda ao número flutuante mais
  próximo; em caso de empate, arredonda ao par
\item
  \textbf{Round toward zero (truncation):} descarta os dígitos fora de
  \(p\) bits
\item
  \textbf{Round toward +infinity:} sempre arredonda para cima
\item
  \textbf{Round toward -infinity:} sempre arredonda para baixo
\end{enumerate}

\subsection{Erro de arredondamento}\label{erro-de-arredondamento}

Se \(x\) é um número real e \(\text{fl}(x)\) é sua representação em
máquina, o erro relativo é limitado por:

\[\frac{|x - \text{fl}(x)|}{|x|} \leq \frac{\epsilon}{2}\]

onde \(\epsilon\) é a machine epsilon.

\begin{center}\rule{0.5\linewidth}{0.5pt}\end{center}

\section{Operações Aritméticas em Ponto
Flutuante}\label{operauxe7uxf5es-aritmuxe9ticas-em-ponto-flutuante}

\subsection{Perda de significância (catastrophic
cancellation)}\label{perda-de-significuxe2ncia-catastrophic-cancellation}

Um problema grave ocorre quando \textbf{subtraímos dois números muito
próximos}.

\begin{tcolorbox}[enhanced jigsaw, arc=.35mm, bottomrule=.15mm, bottomtitle=1mm, breakable, colback=white, colbacktitle=quarto-callout-warning-color!10!white, colframe=quarto-callout-warning-color-frame, coltitle=black, left=2mm, leftrule=.75mm, opacityback=0, opacitybacktitle=0.6, rightrule=.15mm, title=\textcolor{quarto-callout-warning-color}{\faExclamationTriangle}\hspace{0.5em}{Exemplo clássico}, titlerule=0mm, toprule=.15mm, toptitle=1mm]

Calcular \(\sqrt{x^2 + 1} - x\) para \(x = 10^8\).

\textbf{Matematicamente}, para \(x\) grande:
\[\sqrt{x^2 + 1} - x \approx \frac{1}{2x}\]

\textbf{Em máquina}, \(\sqrt{x^2 + 1} \approx x\) (pois \(x\) é muito
grande), então:
\[\sqrt{x^2 + 1} - x \approx 0 \quad \text{❌ resultado errado!}\]

\textbf{Solução:} Racionalizar
\[\sqrt{x^2 + 1} - x = \frac{(x^2 + 1) - x^2}{\sqrt{x^2 + 1} + x} = \frac{1}{\sqrt{x^2 + 1} + x}\]

que é numericamente estável. ✓

\end{tcolorbox}

\begin{Shaded}
\begin{Highlighting}[]
\ImportTok{import}\NormalTok{ math}

\NormalTok{x }\OperatorTok{=} \FloatTok{1e8}

\CommentTok{\# Método instável}
\NormalTok{result1 }\OperatorTok{=}\NormalTok{ math.sqrt(x}\OperatorTok{**}\DecValTok{2} \OperatorTok{+} \DecValTok{1}\NormalTok{) }\OperatorTok{{-}}\NormalTok{ x}
\BuiltInTok{print}\NormalTok{(}\SpecialStringTok{f"Instável: sqrt(x\^{}2 + 1) {-} x = }\SpecialCharTok{\{}\NormalTok{result1}\SpecialCharTok{\}}\SpecialStringTok{"}\NormalTok{)}

\CommentTok{\# Método estável}
\NormalTok{result2 }\OperatorTok{=} \DecValTok{1} \OperatorTok{/}\NormalTok{ (math.sqrt(x}\OperatorTok{**}\DecValTok{2} \OperatorTok{+} \DecValTok{1}\NormalTok{) }\OperatorTok{+}\NormalTok{ x)}
\BuiltInTok{print}\NormalTok{(}\SpecialStringTok{f"Estável:  1/(sqrt(x\^{}2 + 1) + x) = }\SpecialCharTok{\{}\NormalTok{result2}\SpecialCharTok{\}}\SpecialStringTok{"}\NormalTok{)}

\CommentTok{\# Comparação}
\BuiltInTok{print}\NormalTok{(}\SpecialStringTok{f"}\CharTok{\textbackslash{}n}\SpecialStringTok{Valor teórico: 1/(2x) = }\SpecialCharTok{\{}\DecValTok{1}\OperatorTok{/}\NormalTok{(}\DecValTok{2}\OperatorTok{*}\NormalTok{x)}\SpecialCharTok{\}}\SpecialStringTok{"}\NormalTok{)}
\BuiltInTok{print}\NormalTok{(}\SpecialStringTok{f"Erro instável:  }\SpecialCharTok{\{}\BuiltInTok{abs}\NormalTok{(result1 }\OperatorTok{{-}} \DecValTok{1}\OperatorTok{/}\NormalTok{(}\DecValTok{2}\OperatorTok{*}\NormalTok{x))}\SpecialCharTok{\}}\SpecialStringTok{"}\NormalTok{)}
\BuiltInTok{print}\NormalTok{(}\SpecialStringTok{f"Erro estável:   }\SpecialCharTok{\{}\BuiltInTok{abs}\NormalTok{(result2 }\OperatorTok{{-}} \DecValTok{1}\OperatorTok{/}\NormalTok{(}\DecValTok{2}\OperatorTok{*}\NormalTok{x))}\SpecialCharTok{\}}\SpecialStringTok{"}\NormalTok{)}
\end{Highlighting}
\end{Shaded}

\begin{verbatim}
Instável: sqrt(x^2 + 1) - x = 0.0
Estável:  1/(sqrt(x^2 + 1) + x) = 5e-09

Valor teórico: 1/(2x) = 5e-09
Erro instável:  5e-09
Erro estável:   0.0
\end{verbatim}

\subsection{Associatividade quebrada}\label{associatividade-quebrada}

Operações não são associativas em ponto flutuante:
\[(a + b) + c \neq a + (b + c)\]

\begin{Shaded}
\begin{Highlighting}[]
\CommentTok{\# Demonstração}
\NormalTok{a }\OperatorTok{=} \FloatTok{1.0}
\NormalTok{b }\OperatorTok{=} \FloatTok{1e{-}16}
\NormalTok{c }\OperatorTok{=} \OperatorTok{{-}}\FloatTok{1e{-}16}

\NormalTok{result1 }\OperatorTok{=}\NormalTok{ (a }\OperatorTok{+}\NormalTok{ b) }\OperatorTok{+}\NormalTok{ c}
\NormalTok{result2 }\OperatorTok{=}\NormalTok{ a }\OperatorTok{+}\NormalTok{ (b }\OperatorTok{+}\NormalTok{ c)}

\BuiltInTok{print}\NormalTok{(}\SpecialStringTok{f"(a + b) + c = }\SpecialCharTok{\{}\NormalTok{result1}\SpecialCharTok{\}}\SpecialStringTok{"}\NormalTok{)}
\BuiltInTok{print}\NormalTok{(}\SpecialStringTok{f"a + (b + c) = }\SpecialCharTok{\{}\NormalTok{result2}\SpecialCharTok{\}}\SpecialStringTok{"}\NormalTok{)}
\BuiltInTok{print}\NormalTok{(}\SpecialStringTok{f"Iguais? }\SpecialCharTok{\{}\NormalTok{result1 }\OperatorTok{==}\NormalTok{ result2}\SpecialCharTok{\}}\SpecialStringTok{"}\NormalTok{)}

\CommentTok{\# Outro exemplo}
\NormalTok{x }\OperatorTok{=}\NormalTok{ [}\FloatTok{1e10}\NormalTok{, }\DecValTok{1}\NormalTok{, }\OperatorTok{{-}}\FloatTok{1e10}\NormalTok{]}
\BuiltInTok{print}\NormalTok{(}\SpecialStringTok{f"}\CharTok{\textbackslash{}n}\SpecialStringTok{Soma esquerda: ((1e10 + 1) {-} 1e10) = }\SpecialCharTok{\{}\NormalTok{(x[}\DecValTok{0}\NormalTok{] }\OperatorTok{+}\NormalTok{ x[}\DecValTok{1}\NormalTok{]) }\OperatorTok{+}\NormalTok{ x[}\DecValTok{2}\NormalTok{]}\SpecialCharTok{\}}\SpecialStringTok{"}\NormalTok{)}
\BuiltInTok{print}\NormalTok{(}\SpecialStringTok{f"Soma direita:  (1e10 + (1 {-} 1e10)) = }\SpecialCharTok{\{}\NormalTok{x[}\DecValTok{0}\NormalTok{] }\OperatorTok{+}\NormalTok{ (x[}\DecValTok{1}\NormalTok{] }\OperatorTok{+}\NormalTok{ x[}\DecValTok{2}\NormalTok{])}\SpecialCharTok{\}}\SpecialStringTok{"}\NormalTok{)}
\end{Highlighting}
\end{Shaded}

\begin{verbatim}
(a + b) + c = 0.9999999999999999
a + (b + c) = 1.0
Iguais? False

Soma esquerda: ((1e10 + 1) - 1e10) = 1.0
Soma direita:  (1e10 + (1 - 1e10)) = 1.0
\end{verbatim}

\subsection{Absorção numérica}\label{absoruxe7uxe3o-numuxe9rica}

Se \(|a| \gg |b|\), então \(a + b \approx a\) em máquina (o \(b\) é
``absorvido'').

\begin{Shaded}
\begin{Highlighting}[]
\CommentTok{\# Demonstração}
\NormalTok{a }\OperatorTok{=} \FloatTok{1e16}
\NormalTok{b }\OperatorTok{=} \DecValTok{1}

\NormalTok{result }\OperatorTok{=}\NormalTok{ a }\OperatorTok{+}\NormalTok{ b }\OperatorTok{{-}}\NormalTok{ a}
\BuiltInTok{print}\NormalTok{(}\SpecialStringTok{f"(1e16 + 1) {-} 1e16 = }\SpecialCharTok{\{}\NormalTok{result}\SpecialCharTok{\}}\SpecialStringTok{"}\NormalTok{)}
\BuiltInTok{print}\NormalTok{(}\SpecialStringTok{f"Esperado: 1"}\NormalTok{)}
\BuiltInTok{print}\NormalTok{(}\SpecialStringTok{f"Obtido: }\SpecialCharTok{\{}\NormalTok{result}\SpecialCharTok{\}}\SpecialStringTok{"}\NormalTok{)}

\BuiltInTok{print}\NormalTok{(}\SpecialStringTok{f"}\CharTok{\textbackslash{}n}\SpecialStringTok{Comparação:"}\NormalTok{)}
\BuiltInTok{print}\NormalTok{(}\SpecialStringTok{f"1e16 + 1 = }\SpecialCharTok{\{}\NormalTok{a }\OperatorTok{+}\NormalTok{ b}\SpecialCharTok{\}}\SpecialStringTok{"}\NormalTok{)}
\BuiltInTok{print}\NormalTok{(}\SpecialStringTok{f"1e16     = }\SpecialCharTok{\{}\NormalTok{a}\SpecialCharTok{\}}\SpecialStringTok{"}\NormalTok{)}
\BuiltInTok{print}\NormalTok{(}\SpecialStringTok{f"São iguais? }\SpecialCharTok{\{}\NormalTok{a }\OperatorTok{+}\NormalTok{ b }\OperatorTok{==}\NormalTok{ a}\SpecialCharTok{\}}\SpecialStringTok{"}\NormalTok{)}
\end{Highlighting}
\end{Shaded}

\begin{verbatim}
(1e16 + 1) - 1e16 = 0.0
Esperado: 1
Obtido: 0.0

Comparação:
1e16 + 1 = 1e+16
1e16     = 1e+16
São iguais? True
\end{verbatim}

\begin{center}\rule{0.5\linewidth}{0.5pt}\end{center}

\section{Implementação em
Computador}\label{implementauxe7uxe3o-em-computador}

\subsection{Verificação de machine epsilon em
Python}\label{verificauxe7uxe3o-de-machine-epsilon-em-python}

\begin{Shaded}
\begin{Highlighting}[]
\ImportTok{import}\NormalTok{ numpy }\ImportTok{as}\NormalTok{ np}

\CommentTok{\# Método 1: usando numpy}
\NormalTok{eps64 }\OperatorTok{=}\NormalTok{ np.finfo(np.float64).eps}
\BuiltInTok{print}\NormalTok{(}\SpecialStringTok{f"NumPy epsilon (float64): }\SpecialCharTok{\{}\NormalTok{eps64}\SpecialCharTok{\}}\SpecialStringTok{"}\NormalTok{)}

\CommentTok{\# Método 2: encontrar experimentalmente}
\KeywordTok{def}\NormalTok{ find\_machine\_epsilon():}
\NormalTok{    epsilon }\OperatorTok{=} \FloatTok{1.0}
    \ControlFlowTok{while} \FloatTok{1.0} \OperatorTok{+}\NormalTok{ epsilon }\OperatorTok{!=} \FloatTok{1.0}\NormalTok{:}
\NormalTok{        epsilon }\OperatorTok{=}\NormalTok{ epsilon }\OperatorTok{/} \FloatTok{2.0}
    \ControlFlowTok{return}\NormalTok{ epsilon }\OperatorTok{*} \DecValTok{2}

\NormalTok{eps\_experimental }\OperatorTok{=}\NormalTok{ find\_machine\_epsilon()}
\BuiltInTok{print}\NormalTok{(}\SpecialStringTok{f"Experimental:           }\SpecialCharTok{\{}\NormalTok{eps\_experimental}\SpecialCharTok{\}}\SpecialStringTok{"}\NormalTok{)}

\CommentTok{\# Método 3: usando sys}
\ImportTok{import}\NormalTok{ sys}
\BuiltInTok{print}\NormalTok{(}\SpecialStringTok{f"sys.float\_info.epsilon: }\SpecialCharTok{\{}\NormalTok{sys}\SpecialCharTok{.}\NormalTok{float\_info}\SpecialCharTok{.}\NormalTok{epsilon}\SpecialCharTok{\}}\SpecialStringTok{"}\NormalTok{)}
\end{Highlighting}
\end{Shaded}

\begin{verbatim}
NumPy epsilon (float64): 2.220446049250313e-16
Experimental:           2.220446049250313e-16
sys.float_info.epsilon: 2.220446049250313e-16
\end{verbatim}

\subsection{Manipulação de representação
binária}\label{manipulauxe7uxe3o-de-representauxe7uxe3o-binuxe1ria}

\begin{Shaded}
\begin{Highlighting}[]
\KeywordTok{def}\NormalTok{ float\_to\_binary(f):}
    \CommentTok{"""Converte float64 para representação binária"""}
\NormalTok{    bits }\OperatorTok{=}\NormalTok{ struct.pack(}\StringTok{\textquotesingle{}\textgreater{}d\textquotesingle{}}\NormalTok{, f)}
    \ControlFlowTok{return} \StringTok{\textquotesingle{}\textquotesingle{}}\NormalTok{.join(}\SpecialStringTok{f\textquotesingle{}}\SpecialCharTok{\{}\NormalTok{b}\SpecialCharTok{:08b\}}\SpecialStringTok{\textquotesingle{}} \ControlFlowTok{for}\NormalTok{ b }\KeywordTok{in}\NormalTok{ bits)}

\KeywordTok{def}\NormalTok{ analyze\_float(x):}
    \CommentTok{"""Analisa componentes de um float64"""}
\NormalTok{    bits }\OperatorTok{=}\NormalTok{ float\_to\_binary(x)}
\NormalTok{    sign }\OperatorTok{=}\NormalTok{ bits[}\DecValTok{0}\NormalTok{]}
\NormalTok{    exponent }\OperatorTok{=}\NormalTok{ bits[}\DecValTok{1}\NormalTok{:}\DecValTok{12}\NormalTok{]}
\NormalTok{    mantissa }\OperatorTok{=}\NormalTok{ bits[}\DecValTok{12}\NormalTok{:]}
    
\NormalTok{    exp\_value }\OperatorTok{=} \BuiltInTok{int}\NormalTok{(exponent, }\DecValTok{2}\NormalTok{)}
\NormalTok{    exp\_real }\OperatorTok{=}\NormalTok{ exp\_value }\OperatorTok{{-}} \DecValTok{1023}
    
    \BuiltInTok{print}\NormalTok{(}\SpecialStringTok{f"Número: }\SpecialCharTok{\{}\NormalTok{x}\SpecialCharTok{\}}\SpecialStringTok{"}\NormalTok{)}
    \BuiltInTok{print}\NormalTok{(}\SpecialStringTok{f"Binário: }\SpecialCharTok{\{}\NormalTok{bits}\SpecialCharTok{\}}\SpecialStringTok{"}\NormalTok{)}
    \BuiltInTok{print}\NormalTok{(}\SpecialStringTok{f"  Sinal:    }\SpecialCharTok{\{}\NormalTok{sign}\SpecialCharTok{\}}\SpecialStringTok{ (}\SpecialCharTok{\{}\StringTok{\textquotesingle{}+\textquotesingle{}} \ControlFlowTok{if}\NormalTok{ sign }\OperatorTok{==} \StringTok{\textquotesingle{}0\textquotesingle{}} \ControlFlowTok{else} \StringTok{\textquotesingle{}{-}\textquotesingle{}}\SpecialCharTok{\}}\SpecialStringTok{)"}\NormalTok{)}
    \BuiltInTok{print}\NormalTok{(}\SpecialStringTok{f"  Expoente: }\SpecialCharTok{\{}\NormalTok{exponent}\SpecialCharTok{\}}\SpecialStringTok{ (}\SpecialCharTok{\{}\NormalTok{exp\_value}\SpecialCharTok{\}}\SpecialStringTok{ → real: }\SpecialCharTok{\{}\NormalTok{exp\_real}\SpecialCharTok{\}}\SpecialStringTok{)"}\NormalTok{)}
    \BuiltInTok{print}\NormalTok{(}\SpecialStringTok{f"  Mantissa: }\SpecialCharTok{\{}\NormalTok{mantissa[:}\DecValTok{16}\NormalTok{]}\SpecialCharTok{\}}\SpecialStringTok{... (}\SpecialCharTok{\{}\NormalTok{mantissa}\SpecialCharTok{\}}\SpecialStringTok{)"}\NormalTok{)}
    \BuiltInTok{print}\NormalTok{()}

\CommentTok{\# Exemplos}
\NormalTok{analyze\_float(}\FloatTok{1.0}\NormalTok{)}
\NormalTok{analyze\_float(}\FloatTok{0.5}\NormalTok{)}
\NormalTok{analyze\_float(}\FloatTok{0.1}\NormalTok{)}
\end{Highlighting}
\end{Shaded}

\begin{verbatim}
Número: 1.0
Binário: 0011111111110000000000000000000000000000000000000000000000000000
  Sinal:    0 (+)
  Expoente: 01111111111 (1023 → real: 0)
  Mantissa: 0000000000000000... (0000000000000000000000000000000000000000000000000000)

Número: 0.5
Binário: 0011111111100000000000000000000000000000000000000000000000000000
  Sinal:    0 (+)
  Expoente: 01111111110 (1022 → real: -1)
  Mantissa: 0000000000000000... (0000000000000000000000000000000000000000000000000000)

Número: 0.1
Binário: 0011111110111001100110011001100110011001100110011001100110011010
  Sinal:    0 (+)
  Expoente: 01111111011 (1019 → real: -4)
  Mantissa: 1001100110011001... (1001100110011001100110011001100110011001100110011010)
\end{verbatim}

\subsection{Verificação de
overflow/underflow}\label{verificauxe7uxe3o-de-overflowunderflow}

\begin{Shaded}
\begin{Highlighting}[]
\ImportTok{import}\NormalTok{ numpy }\ImportTok{as}\NormalTok{ np}

\CommentTok{\# Limites}
\NormalTok{info }\OperatorTok{=}\NormalTok{ np.finfo(np.float64)}
\BuiltInTok{print}\NormalTok{(}\StringTok{"Float64 limits:"}\NormalTok{)}
\BuiltInTok{print}\NormalTok{(}\SpecialStringTok{f"  Epsilon:            }\SpecialCharTok{\{}\NormalTok{info}\SpecialCharTok{.}\NormalTok{eps}\SpecialCharTok{\}}\SpecialStringTok{"}\NormalTok{)}
\BuiltInTok{print}\NormalTok{(}\SpecialStringTok{f"  Tiny (min normal):  }\SpecialCharTok{\{}\NormalTok{info}\SpecialCharTok{.}\NormalTok{tiny}\SpecialCharTok{\}}\SpecialStringTok{"}\NormalTok{)}
\BuiltInTok{print}\NormalTok{(}\SpecialStringTok{f"  Max:                }\SpecialCharTok{\{}\NormalTok{info}\SpecialCharTok{.}\BuiltInTok{max}\SpecialCharTok{\}}\SpecialStringTok{"}\NormalTok{)}

\CommentTok{\# Gerar overflow}
\NormalTok{x }\OperatorTok{=}\NormalTok{ np.float64(}\FloatTok{1e308}\NormalTok{)}
\BuiltInTok{print}\NormalTok{(}\SpecialStringTok{f"}\CharTok{\textbackslash{}n}\SpecialStringTok{1e308 * 10 = }\SpecialCharTok{\{}\NormalTok{x }\OperatorTok{*} \DecValTok{10}\SpecialCharTok{\}}\SpecialStringTok{"}\NormalTok{)  }\CommentTok{\# inf}

\CommentTok{\# Underflow gradual}
\NormalTok{y }\OperatorTok{=}\NormalTok{ np.float64(}\FloatTok{1e{-}320}\NormalTok{)}
\BuiltInTok{print}\NormalTok{(}\SpecialStringTok{f"1e{-}320 = }\SpecialCharTok{\{}\NormalTok{y}\SpecialCharTok{\}}\SpecialStringTok{"}\NormalTok{)}
\BuiltInTok{print}\NormalTok{(}\SpecialStringTok{f"1e{-}330 = }\SpecialCharTok{\{}\NormalTok{np}\SpecialCharTok{.}\NormalTok{float64(}\FloatTok{1e{-}330}\NormalTok{)}\SpecialCharTok{\}}\SpecialStringTok{"}\NormalTok{)  }\CommentTok{\# Desnormalizado}
\end{Highlighting}
\end{Shaded}

\begin{verbatim}
Float64 limits:
  Epsilon:            2.220446049250313e-16
  Tiny (min normal):  2.2250738585072014e-308
  Max:                1.7976931348623157e+308

1e308 * 10 = inf
1e-320 = 1e-320
1e-330 = 0.0
\end{verbatim}

\begin{verbatim}
/var/folders/bm/kxnd_r6x0vd3_b38fvf0rrgr0000gn/T/ipykernel_11751/740743619.py:12: RuntimeWarning: overflow encountered in scalar multiply
  print(f"\n1e308 * 10 = {x * 10}")  # inf
\end{verbatim}

\begin{center}\rule{0.5\linewidth}{0.5pt}\end{center}

\section{Aplicações e Conexões com Cálculo
Numérico}\label{aplicauxe7uxf5es-e-conexuxf5es-com-cuxe1lculo-numuxe9rico}

\subsection{Por que isso importa}\label{por-que-isso-importa}

\begin{enumerate}
\def\labelenumi{\arabic{enumi}.}
\tightlist
\item
  \textbf{Propagação de erros:} Pequenos erros de arredondamento crescem
  ao longo de algoritmos iterativos
\item
  \textbf{Escolha de algoritmo:} Métodos numericamente estáveis são
  essenciais
\item
  \textbf{Análise de erro:} Devemos estimar como erros se propagam em
  nossas contas
\item
  \textbf{Teste de convergência:} Saber \(\epsilon\) ajuda a definir
  tolerâncias apropriadas
\end{enumerate}

\subsection{Próximos capítulos}\label{pruxf3ximos-capuxedtulos}

\begin{itemize}
\tightlist
\item
  \textbf{Cap 1: Propagação de Erros} --- Como estimar e controlar erros
\item
  \textbf{Cap 2: Sistemas Lineares} --- Resolução estável via métodos
  diretos e iterativos
\item
  \textbf{Cap 3: Interpolação} --- Aproximação de funções e análise de
  erro
\item
  \textbf{Cap 4: Integração Numérica} --- Quadraturas e fórmulas
  compostas
\end{itemize}

\begin{center}\rule{0.5\linewidth}{0.5pt}\end{center}

\section{Exercícios Propostos}\label{exercuxedcios-propostos}

\begin{enumerate}
\def\labelenumi{\arabic{enumi}.}
\tightlist
\item
  \textbf{Conversão de bases}

  \begin{itemize}
  \tightlist
  \item
    Converter \((101.11)_2\) para decimal
  \item
    Converter \((0.1)_{10}\) para binário (até 16 dígitos)
  \end{itemize}
\item
  \textbf{Machine epsilon}

  \begin{itemize}
  \tightlist
  \item
    Qual é o machine epsilon para float32 e float64?
  \item
    Implementar um algoritmo em Python que encontre experimentalmente
  \end{itemize}
\item
  \textbf{Limites numéricos}

  \begin{itemize}
  \tightlist
  \item
    Em double precision IEEE 754, qual é o maior expoente?
  \item
    Qual é o valor máximo representável?
  \end{itemize}
\item
  \textbf{Não-associatividade}

  \begin{itemize}
  \tightlist
  \item
    Verificar numericamente que \((a + b) + c \neq a + (b + c)\) com
    valores apropriados
  \end{itemize}
\item
  \textbf{Cancelamento numérico}

  \begin{itemize}
  \tightlist
  \item
    Implementar ambas as formulações de \(\sqrt{x^2 + 1} - x\) em Python
  \item
    Comparar resultados para \(x = 10^{10}\)
  \end{itemize}
\item
  \textbf{Absorção numérica}

  \begin{itemize}
  \tightlist
  \item
    Verificar que \((10^{16} + 1) - 10^{16} = 0\) em float64
  \item
    Explorar diferentes valores de \(10^k\) para encontrar threshold
  \end{itemize}
\end{enumerate}

\begin{center}\rule{0.5\linewidth}{0.5pt}\end{center}

\section{Referências}\label{referuxeancias}

\begin{itemize}
\item
  \textbf{Burden, R. L., \& Faires, J. D. (2016).} \emph{Numerical
  analysis} (10th ed.). Cengage Learning.
\item
  \textbf{Goldberg, D. (1991).} What every computer scientist should
  know about floating-point arithmetic. \emph{ACM Computing Surveys},
  23(1), 5--48.
\item
  \textbf{Higham, N. J. (2002).} \emph{Accuracy and stability of
  numerical algorithms} (2nd ed.). SIAM.
\item
  \textbf{IEEE. (2019).} IEEE 754-2019: IEEE standard for floating-point
  arithmetic. Institute of Electrical and Electronics Engineers.
\item
  \textbf{Knuth, D. E. (1997).} \emph{The art of computer programming:
  Seminumerical algorithms} (Vol. 2, 3rd ed.). Addison-Wesley.
\end{itemize}

\begin{center}\rule{0.5\linewidth}{0.5pt}\end{center}

\section{Anexo: Código Python Completo para
Experimentos}\label{anexo-cuxf3digo-python-completo-para-experimentos}

\begin{Shaded}
\begin{Highlighting}[]
\CommentTok{\# Código para experimentos com representação numérica}

\ImportTok{import}\NormalTok{ numpy }\ImportTok{as}\NormalTok{ np}
\ImportTok{import}\NormalTok{ struct}
\ImportTok{import}\NormalTok{ sys}

\KeywordTok{class}\NormalTok{ FloatAnalyzer:}
    \CommentTok{"""Classe para analisar representações de ponto flutuante"""}
    
    \AttributeTok{@staticmethod}
    \KeywordTok{def}\NormalTok{ to\_binary(x):}
        \CommentTok{"""Converte float64 para binário"""}
\NormalTok{        bits }\OperatorTok{=}\NormalTok{ struct.pack(}\StringTok{\textquotesingle{}\textgreater{}d\textquotesingle{}}\NormalTok{, x)}
        \ControlFlowTok{return} \StringTok{\textquotesingle{}\textquotesingle{}}\NormalTok{.join(}\SpecialStringTok{f\textquotesingle{}}\SpecialCharTok{\{}\NormalTok{b}\SpecialCharTok{:08b\}}\SpecialStringTok{\textquotesingle{}} \ControlFlowTok{for}\NormalTok{ b }\KeywordTok{in}\NormalTok{ bits)}
    
    \AttributeTok{@staticmethod}
    \KeywordTok{def}\NormalTok{ from\_binary(bits\_str):}
        \CommentTok{"""Converte binário para float64"""}
\NormalTok{        bytes\_val }\OperatorTok{=} \BuiltInTok{bytes}\NormalTok{(}\BuiltInTok{int}\NormalTok{(bits\_str[i:i}\OperatorTok{+}\DecValTok{8}\NormalTok{], }\DecValTok{2}\NormalTok{) }
                         \ControlFlowTok{for}\NormalTok{ i }\KeywordTok{in} \BuiltInTok{range}\NormalTok{(}\DecValTok{0}\NormalTok{, }\DecValTok{64}\NormalTok{, }\DecValTok{8}\NormalTok{))}
        \ControlFlowTok{return}\NormalTok{ struct.unpack(}\StringTok{\textquotesingle{}\textgreater{}d\textquotesingle{}}\NormalTok{, bytes\_val)[}\DecValTok{0}\NormalTok{]}
    
    \AttributeTok{@staticmethod}
    \KeywordTok{def}\NormalTok{ analyze(x):}
        \CommentTok{"""Analisa componentes de float64"""}
\NormalTok{        bits }\OperatorTok{=}\NormalTok{ FloatAnalyzer.to\_binary(x)}
\NormalTok{        sign\_bit }\OperatorTok{=}\NormalTok{ bits[}\DecValTok{0}\NormalTok{]}
\NormalTok{        exp\_bits }\OperatorTok{=}\NormalTok{ bits[}\DecValTok{1}\NormalTok{:}\DecValTok{12}\NormalTok{]}
\NormalTok{        mant\_bits }\OperatorTok{=}\NormalTok{ bits[}\DecValTok{12}\NormalTok{:]}
        
\NormalTok{        exp\_val }\OperatorTok{=} \BuiltInTok{int}\NormalTok{(exp\_bits, }\DecValTok{2}\NormalTok{)}
\NormalTok{        exp\_real }\OperatorTok{=}\NormalTok{ exp\_val }\OperatorTok{{-}} \DecValTok{1023}
        
        \ControlFlowTok{return}\NormalTok{ \{}
            \StringTok{\textquotesingle{}valor\textquotesingle{}}\NormalTok{: x,}
            \StringTok{\textquotesingle{}binario\textquotesingle{}}\NormalTok{: bits,}
            \StringTok{\textquotesingle{}sinal\textquotesingle{}}\NormalTok{: sign\_bit,}
            \StringTok{\textquotesingle{}expoente\_bits\textquotesingle{}}\NormalTok{: exp\_bits,}
            \StringTok{\textquotesingle{}expoente\_valor\textquotesingle{}}\NormalTok{: exp\_val,}
            \StringTok{\textquotesingle{}expoente\_real\textquotesingle{}}\NormalTok{: exp\_real,}
            \StringTok{\textquotesingle{}mantissa\textquotesingle{}}\NormalTok{: mant\_bits}
\NormalTok{        \}}

\CommentTok{\# Exemplos de uso}
\BuiltInTok{print}\NormalTok{(}\StringTok{"=== Análise de Números em IEEE 754 ===}\CharTok{\textbackslash{}n}\StringTok{"}\NormalTok{)}

\ControlFlowTok{for}\NormalTok{ num }\KeywordTok{in}\NormalTok{ [}\FloatTok{1.0}\NormalTok{, }\FloatTok{0.5}\NormalTok{, }\FloatTok{2.0}\NormalTok{, }\FloatTok{0.1}\NormalTok{]:}
\NormalTok{    result }\OperatorTok{=}\NormalTok{ FloatAnalyzer.analyze(num)}
    \BuiltInTok{print}\NormalTok{(}\SpecialStringTok{f"Número: }\SpecialCharTok{\{}\NormalTok{num}\SpecialCharTok{\}}\SpecialStringTok{"}\NormalTok{)}
    \BuiltInTok{print}\NormalTok{(}\SpecialStringTok{f"  Expoente real: }\SpecialCharTok{\{}\NormalTok{result[}\StringTok{\textquotesingle{}expoente\_real\textquotesingle{}}\NormalTok{]}\SpecialCharTok{\}}\SpecialStringTok{"}\NormalTok{)}
    \BuiltInTok{print}\NormalTok{(}\SpecialStringTok{f"  Binário: }\SpecialCharTok{\{}\NormalTok{result[}\StringTok{\textquotesingle{}binario\textquotesingle{}}\NormalTok{][:}\DecValTok{16}\NormalTok{]}\SpecialCharTok{\}}\SpecialStringTok{..."}\NormalTok{)}
    \BuiltInTok{print}\NormalTok{()}
\end{Highlighting}
\end{Shaded}

\begin{verbatim}
=== Análise de Números em IEEE 754 ===

Número: 1.0
  Expoente real: 0
  Binário: 0011111111110000...

Número: 0.5
  Expoente real: -1
  Binário: 0011111111100000...

Número: 2.0
  Expoente real: 1
  Binário: 0100000000000000...

Número: 0.1
  Expoente real: -4
  Binário: 0011111110111001...
\end{verbatim}

\begin{center}\rule{0.5\linewidth}{0.5pt}\end{center}

\emph{Última atualização: 2024}




\end{document}
